% GENERATED FILE. Edit canonical lessons, then run npm run generate:books.
% canonical-source-hash: fnv1a64:4e3cd35c75e79118
% canonical-lessons: LA-C397-iusiurandum, LA-C397-conventio, LA-C397-paciscor, LA-C397-pactum, LA-C397-stipulor

\chapter{Oaths and Agreements}
\label{ch:la-oaths-and-agreements}
\hlchaptermodality{397}

\begin{chapteropening}
\textbf{By the end of this chapter:} \emph{I can talk about an oath, an agreement and a pact, and say that people strike a bargain and demand a formal promise.}

Complete the last lesson of chapter 397: I can talk about an oath, an agreement and a pact, and say that people strike a bargain and demand a formal promise.
\end{chapteropening}

\section[iūsiūrandum, iūrisiūrandī]{iūsiūrandum, iūrisiūrandī — an oath}

\label{lesson:LA-C397-iusiurandum}



\begin{quote}
Before the new one: say the Latin for a solemn promise, then the Latin for to bind.
\end{quote}

\subsection*{You'll want to know: iūsiūrandum, iūrisiūrandī}
\textbf{iūsiūrandum, iūrisiūrandī} — \textquotedblleft{}an oath\textquotedblright{}.

\begin{quote}
Iūsiūrandum dedit sē reditūrum esse.
\end{quote}

\begin{quote}
\emph{He gave his oath that he would return.}
\end{quote}

\emph{iūsiūrandum, iūrisiūrandī} n.: both parts decline, \emph{iūre iūrandō}, by oath. Built from \emph{iūs} and \emph{iūrō}, to swear.

\begin{cousinweb}
\emph{iūsiūrandum} is built from \emph{iūs} \textquotedblleft{}right, law\textquotedblright{} and \emph{iūrandum}, the gerundive of \emph{iūrō} \textquotedblleft{}to swear\textquotedblright{} (Lewis \& Short): a right that must be sworn.
\end{cousinweb}

\begin{grammarlens}[title={What you've built}]
One more word for telling, explaining and describing at B1.
\end{grammarlens}

\subsection*{Your turn}
\begin{itemize}
\raggedright
  \item \emph{Say it:} \emph{iūsiūrandum, iūrisiūrandī}
  \item \emph{Say it:} \emph{iūsiūrandum, iūrisiūrandī}, once more
  \item \emph{Say it:} a sentence of your own with \emph{iūsiūrandum, iūrisiūrandī}
  \item \emph{Recall:} say the Latin for a solemn promise, then the Latin for to bind, then say \emph{iūsiūrandum, iūrisiūrandī} again
\end{itemize}

\begin{tcolorbox}[breakable,colback=teal!4,colframe=teal!35!black,title={Before you move on}]
What does iūsiūrandum, iūrisiūrandī mean? (\textquotedblleft{}An oath\textquotedblright{}.) Say it once more.
\end{tcolorbox}
\section[cōnventiō, cōnventiōnis]{cōnventiō, cōnventiōnis — an agreement, a covenant; an assembly}

\label{lesson:LA-C397-conventio}



\begin{quote}
Before the new one: say the Latin for to bind, then the Latin for an oath.
\end{quote}

\subsection*{You'll want to know: cōnventiō, cōnventiōnis}
\textbf{cōnventiō, cōnventiōnis} — \textquotedblleft{}an agreement, a covenant; an assembly\textquotedblright{}.

\begin{quote}
Cōnventiō inter eōs facta est.
\end{quote}

\begin{quote}
\emph{An agreement was made between them.}
\end{quote}

\emph{cōnventiō, cōnventiōnis} f., built on \emph{conveniō}, to come together: \emph{convenit}, it is agreed. English \emph{convention}.

\begin{cousinweb}
\emph{cōnventiō} is the noun built on \emph{conveniō} \textquotedblleft{}to come together, to agree\textquotedblright{} (Lewis \& Short). English \emph{convention} comes from it.
\end{cousinweb}

\begin{grammarlens}[title={What you've built}]
One more word for telling, explaining and describing at B1.
\end{grammarlens}

\subsection*{Your turn}
\begin{itemize}
\raggedright
  \item \emph{Say it:} \emph{cōnventiō, cōnventiōnis}
  \item \emph{Say it:} \emph{cōnventiō, cōnventiōnis}, once more
  \item \emph{Say it:} a sentence of your own with \emph{cōnventiō, cōnventiōnis}
  \item \emph{Recall:} say the Latin for to bind, then the Latin for an oath, then say \emph{cōnventiō, cōnventiōnis} again
\end{itemize}

\begin{tcolorbox}[breakable,colback=teal!4,colframe=teal!35!black,title={Before you move on}]
What does cōnventiō, cōnventiōnis mean? (\textquotedblleft{}An agreement, a covenant; an assembly\textquotedblright{}.) Say it once more.
\end{tcolorbox}
\section[pacīscor, pacīscī]{pacīscor, pacīscī — to strike a bargain, to come to terms}

\label{lesson:LA-C397-paciscor}



\begin{quote}
Before the new one: say the Latin for an oath, then the Latin for an agreement.
\end{quote}

\subsection*{You'll want to know: pacīscor, pacīscī}
\textbf{pacīscor, pacīscī} — \textquotedblleft{}to strike a bargain, to come to terms\textquotedblright{}.

\begin{quote}
Cum hostibus dē pāce pacīscēbantur.
\end{quote}

\begin{quote}
\emph{They were bargaining with the enemy about peace.}
\end{quote}

Deponent: \emph{pacīscor, pacīscī, pactus sum}, kin to \emph{pāx}, peace. \emph{Cum aliquō pacīscī}, to strike a bargain with someone.

\begin{cousinweb}
\emph{pacīscor} is built on an old verb \emph{pacō} \textquotedblleft{}to come to terms\textquotedblright{} with the ending \emph{-scō} (EtymDB; Lewis \& Short), and belongs with \emph{pāx} \textquotedblleft{}peace\textquotedblright{}: to fix terms.
\end{cousinweb}

\begin{grammarlens}[title={What you've built}]
One more word for telling, explaining and describing at B1.
\end{grammarlens}

\subsection*{Your turn}
\begin{itemize}
\raggedright
  \item \emph{Say it:} \emph{pacīscor, pacīscī}
  \item \emph{Say it:} \emph{pacīscor, pacīscī}, once more
  \item \emph{Say it:} a sentence of your own with \emph{pacīscor, pacīscī}
  \item \emph{Recall:} say the Latin for an oath, then the Latin for an agreement, then say \emph{pacīscor, pacīscī} again
\end{itemize}

\begin{tcolorbox}[breakable,colback=teal!4,colframe=teal!35!black,title={Before you move on}]
What does pacīscor, pacīscī mean? (\textquotedblleft{}To strike a bargain, to come to terms\textquotedblright{}.) Say it once more.
\end{tcolorbox}
\section[pactum, pactī]{pactum, pactī — a pact, an agreement}

\label{lesson:LA-C397-pactum}



\begin{quote}
Before the new one: say the Latin for an agreement, then the Latin for to strike a bargain.
\end{quote}

\subsection*{You'll want to know: pactum, pactī}
\textbf{pactum, pactī} — \textquotedblleft{}a pact, an agreement\textquotedblright{}.

\begin{quote}
Pactum servāre dēbēmus.
\end{quote}

\begin{quote}
\emph{We must keep the pact.}
\end{quote}

\emph{pactum, pactī} n., the participle of \emph{pacīscor}: a thing agreed. \emph{Quō pactō?}, in what way? English \emph{pact}.

\begin{cousinweb}
\emph{pactum} is the participle of \emph{pacīscor} \textquotedblleft{}to make a bargain\textquotedblright{} (EtymDB; Lewis \& Short): a thing agreed. English \emph{pact} comes from it.
\end{cousinweb}

\begin{grammarlens}[title={What you've built}]
One more word for telling, explaining and describing at B1.
\end{grammarlens}

\subsection*{Your turn}
\begin{itemize}
\raggedright
  \item \emph{Say it:} \emph{pactum, pactī}
  \item \emph{Say it:} \emph{pactum, pactī}, once more
  \item \emph{Say it:} a sentence of your own with \emph{pactum, pactī}
  \item \emph{Recall:} say the Latin for an agreement, then the Latin for to strike a bargain, then say \emph{pactum, pactī} again
\end{itemize}

\begin{tcolorbox}[breakable,colback=teal!4,colframe=teal!35!black,title={Before you move on}]
What does pactum, pactī mean? (\textquotedblleft{}A pact, an agreement\textquotedblright{}.) Say it once more.
\end{tcolorbox}
\section[stipulor, stipulārī]{stipulor, stipulārī — to bargain for, to demand a formal promise}

\label{lesson:LA-C397-stipulor}



\begin{quote}
Before the new one: say the Latin for to strike a bargain, then the Latin for a pact.
\end{quote}

\subsection*{You'll want to know: stipulor, stipulārī}
\textbf{stipulor, stipulārī} — \textquotedblleft{}to bargain for, to demand a formal promise\textquotedblright{}.

\begin{quote}
Stipulātus est ut pecūnia redderētur.
\end{quote}

\begin{quote}
\emph{He stipulated that the money be returned.}
\end{quote}

Deponent: \emph{stipulor, stipulārī, stipulātus sum}. In Roman law the \emph{stipulātiō} was a contract made by question and answer. English \emph{stipulate}.

\begin{cousinweb}
The origin of \emph{stipulor} is not settled: Lewis \& Short give no sure account, and EtymDB has no parent for it. English \emph{stipulate} comes from it.
\end{cousinweb}

\begin{grammarlens}[title={What you've built}]
One more word for telling, explaining and describing at B1.
\end{grammarlens}

\subsection*{Your turn}
\begin{itemize}
\raggedright
  \item \emph{Say it:} \emph{stipulor, stipulārī}
  \item \emph{Say it:} \emph{stipulor, stipulārī}, once more
  \item \emph{Say it:} a sentence of your own with \emph{stipulor, stipulārī}
  \item \emph{Recall:} say the Latin for to strike a bargain, then the Latin for a pact, then say \emph{stipulor, stipulārī} again
\end{itemize}

\begin{tcolorbox}[breakable,colback=teal!4,colframe=teal!35!black,title={Before you move on}]
What does stipulor, stipulārī mean? (\textquotedblleft{}To bargain for, to demand a formal promise\textquotedblright{}.) Say it once more.
\end{tcolorbox}
